\chapter{$L^p$ 空间、收敛定理与乘积测度}
\label{ch:lp-product-measure}

\begin{introduction}
  \item 先修：第 3 章的可测函数、Lebesgue 积分、几乎处处与支配收敛。
  \item 学习产出：能选择合适的 $L^p$ 误差尺度，并在交换极限或积分次序前逐条核对假设。
  \item 研究连接：预测误差、随机变量矩、联合分布和嵌套期望都依赖本章的共同语言。
\end{introduction}

同一个误差函数在均值、均方和最坏情形下会得到不同判断；同一个二元函数在没有绝对可积性时，甚至会因求和或积分次序不同而得到不同答案。本章把这些现象统一到 $L^p$ 空间、收敛定理与乘积测度中。

\section{$L^p$ 范数与几乎处处等价}

在测度空间 $(\Omega,\mathcal F,\mu)$ 上，对 $1\le p<\infty$ 定义
\[
  \lVert f\rVert_p=\left(\int_\Omega |f|^p\,d\mu\right)^{1/p},
\]
并令 $L^p(\mu)$ 包含所有该值有限的可测函数。$L^\infty$ 使用本质上确界：忽略零测集后仍能成立的最小统一上界。

严格地说，$L^p$ 的元素不是单个函数，而是“几乎处处相等”的等价类。若只在零测集上改变 $f$，所有 $L^p$ 范数不变；因此 $\lVert f\rVert_p=0$ 意味着 $f=0$ 几乎处处，而非每一点都为零。商掉这类差异后，$L^p$ 才成为真正的范数空间，并且对 $1\le p\le\infty$ 都是完备空间。

\begin{note}
在概率空间中，若 $1\le p\le q\le\infty$，则 $L^q\subseteq L^p$；在一般无限测度空间上不能照搬这个包含关系。集合总测度是不可省略的假设。
\end{note}

Hölder 不等式在 $1/p+1/q=1$ 时给出
\[
  \int |fg|\,d\mu\le \lVert f\rVert_p\lVert g\rVert_q,
\]
Minkowski 不等式则给出三角不等式 $\lVert f+g\rVert_p\le\lVert f\rVert_p+\lVert g\rVert_p$。前者控制乘积，后者使误差尺度成为范数。

\subsection{固定账本：$f(x)=x$ 的三个范数}

在 $[0,1]$ 上，
\[
  \lVert x\rVert_p=\left(\int_0^1x^p\,dx\right)^{1/p}
  =(p+1)^{-1/p}.
\]
所以 $\lVert x\rVert_1=0.5$、$\lVert x\rVert_2=0.577350\ldots$、$\lVert x\rVert_4=0.668740\ldots$。$p$ 增大时更重视峰值，但有限 $p$ 仍不是最大误差。

\section{三种收敛门禁不能互换}

第 3 章的三条结果现在可用 $L^p$ 语言定位：
\begin{itemize}[nosep]
  \item 若 $0\le f_n\uparrow f$，单调收敛定理给出积分收敛；
  \item 若 $f_n\to f$ 几乎处处且 $|f_n|\le g\in L^1$，支配收敛定理给出 $L^1$ 意义下的积分交换；
  \item 只有非负性时，Fatou 引理通常只给下半连续不等式。
\end{itemize}

几乎处处收敛、依测度收敛和 $L^p$ 收敛是不同概念。$L^p$ 收敛推出依测度收敛；在有限测度空间中，从依测度收敛的序列可抽出几乎处处收敛子列，但原序列未必几乎处处收敛。第 3 章的尖峰 $n\mathbf 1_{(0,1/n]}$ 正说明：仅有逐点收敛而没有共同可积控制，积分可以完全不收敛。

\section{乘积测度与两种积分次序}

给定两个 $\sigma$ 有限测度空间，乘积 $\sigma$ 代数由可测矩形 $A\times B$ 生成，乘积测度满足
\[
  (\mu\times\nu)(A\times B)=\mu(A)\nu(B).
\]
Tonelli 定理处理非负可测函数：允许交换积分次序，结果也可以是 $+\infty$。Fubini 定理处理可积符号函数：若
\[
  \int |f|\,d(\mu\times\nu)<\infty,
\]
则几乎所有截面可积，两种迭代积分都存在且等于乘积空间积分。非负性与绝对可积性分别是两道不同门禁。

\subsection{失败见证：条件求和不能冒充 Fubini}

在 $\mathbb N\times\mathbb N$ 上取计数测度，并令
\[
  a_{mn}=\begin{cases}
    1,&m=n,\\
    -1,&m=n+1,\\
    0,&\text{其他}.
  \end{cases}
\]
固定一行求和时，第 1 行为 1，其余行中 $1$ 与 $-1$ 抵消，所以先逐行求和得到 1。固定一列求和时，每列的 $1$ 都与下一行的 $-1$ 抵消，所以先逐列求和得到 0。

矛盾来自缺少绝对可积性：前 $N\times N$ 方阵的绝对值和为 $2N-1\to\infty$。方阵截断的有符号和恒为 1，而 $(N+1)\times N$ 截断恒为 0；不同扩张路径保留了不同边界项。Fubini 不适用，Tonelli 也不能用于含负项的 $a_{mn}$。

\section{独立数值实验与误差界}

对 $x^2$ 使用 $N=64$ 个等宽小区间的复合中点公式，程序得到积分 $0.333312988$。解析积分是 $1/3$；由于 $\max|(x^2)''|=2$，中点余项满足
\[
  |E_N|\le\frac{2}{24N^2}=0.000020345\ldots,
\]
实际误差也恰为 $0.000020345\ldots$。双重级数部分只做整数计数，因此误差界为零；它不是用浮点近似制造的次序差异。

\begin{lstlisting}[language=bash]
uv run python tools/check_learning_unit.py \
  --manifest curriculum/manifest.json --unit upper.ch04
\end{lstlisting}

\section*{章末练习}
\begingroup\small
\subsection*{口述概念}
\begin{enumerate}[nosep,leftmargin=*]
  \item 为什么 $L^p$ 的元素是 a.e. 等价类？普通上确界与本质上确界有什么差别？
  \item 比较 Hölder 与 Minkowski 的输入和结论；为什么“$X,Y\in L^1$”不足以保证协方差存在？
  \item 区分 $L^p$、依测度和 a.e. 收敛，并陈述本章可以保证的推出关系。
\end{enumerate}

\subsection*{笔试推导}
\begin{enumerate}[nosep,leftmargin=*]
  \item 从 Young 不等式证明 Hölder，再用 Hölder 推导 $1<p<\infty$ 时的 Minkowski。
  \item 若 $\mu(\Omega)<\infty$ 且 $1\leq p<q\leq\infty$，证明 $\lVert f\rVert_p\leq\mu(\Omega)^{1/p-1/q}\lVert f\rVert_q$（约定 $1/\infty=0$）；给出无限测度下两个包含方向各自失败的函数。
  \item 对双重级数 $a_{mn}$，分别计算两种迭代和、$N\times N$ 与 $(N+1)\times N$ 截断和，并证明绝对值和为 $2N-1$。
\end{enumerate}

\subsection*{数值编程}
\begin{enumerate}[leftmargin=*]
  \item 画出 $n^a\mathbf1_{(0,1/n]}$ 的 $L^p$ 范数，比较 $a=1/3,p=1,2,3,4$。
  \item 构造正文正负对角矩阵，比较 $N\times N$ 与 $(N+1)\times N$ 的和。
  \item 改变单个点上的函数值，比较普通上确界与本质上确界。
\end{enumerate}

\subsection*{研究判断}
\begin{enumerate}[nosep,leftmargin=*]
  \item 某模型报告训练损失在样本点逐点下降到零。需要补什么证据才能声称总体 $L^1$ 风险趋零？
  \item 两段代码分别先按资产、先按日期聚合，得到不同风险值。设计一份区分数学换序失败、缺失数据规则与浮点误差的审计。
  \item 研究员在无限时间轴上用“高阶矩控制低阶矩”删去一个风险检查。评价该论证，并给出最小修制方案。
\end{enumerate}
\endgroup

\subsection*{基础衔接：独立计算}
令 $f_n=n^{1/3}\mathbf1_{(0,1/n]}$。分别判断 $L^1,L^2,L^3$ 收敛到零是否成立。

\subsection*{分级提示}
\begingroup\small
\begin{enumerate}[nosep,leftmargin=*]
  \item \textbf{口述概念：}先区分函数与商空间；零测异常不应改变范数。Hölder 把 $L^p\times L^q$ 送入 $L^1$，Minkowski 控制同一 $L^p$ 中的和。用 Markov 不等式画出 $L^p\Rightarrow$ 依测度，再用子列和 typewriter 说明其余箭头。
  \item \textbf{笔试推导：}先归一化 $f,g$ 再逐点使用 Young；Minkowski 对 $|f+g|^{p-1}$ 使用共轭指数。有限测度嵌入对 $|f|^p\cdot1$ 使用 Hölder。双重级数只需数两条对角线。
  \item \textbf{数值编程：}先写出本章公式中的目标量，再显示中间计算。改变参数前预测结果方向，运行后说明差异来自模型、抽样还是数值近似。
  \item \textbf{研究判断：}逐点样本行为不是总体积分结论；寻找统一可积支配、单调性或直接的 $L^1$ 上界。换序审计先保留整数/高精度独立记录，再核对截断区域和缺失规则。无限测度下必须重新证明嵌入或改用局部/加权空间。
\end{enumerate}
\endgroup


\section*{本章小结与 Capstone 连接}

$L^p$ 把误差尺度变成完备空间；收敛定理规定极限交换条件；Tonelli 与 Fubini 规定积分次序门禁。上册 Capstone 将要求每个期望、联合积分和数值近似同时声明空间、可积条件、独立 oracle 与误差界。
